%============================================================
% CH1 · PATTERN 4 — EXERCISES, HINTS, ANSWERS
%============================================================
\subsection{Pattern 4: Exercises and Going Further}
\label{ssec:ch1-p4-ex}

\begin{patternbox}{Pattern 4 Summary Card}
\begin{tabularx}{\linewidth}{@{}>{\bfseries}l X@{}}
Recognise & the statement describes a process; numbers given digit by digit; positions computed by a formula; ``in place'' with dependencies between cells \\
Core loops & apply the rule step by step; carry loop $s = a_i + b_i + c$ from the least significant end; index maps $i \to a[i]$, $2i$, $\lfloor i/2 \rfloor$, $i \bmod n$, $rC + c$ \\
Tools & update order (right to left for a right shift); encoding $\text{old} + n\cdot\text{new}$; sentinels; a counter instead of updating all later cells; a period instead of many steps \\
Cost & (steps) $\times$ (cost per step) --- check it before coding \\
Top traps & reading values that were already overwritten; converting digit arrays to integers; $R$ instead of $C$ in $rC + c$; negative \texttt{\%}; the extra digit of a final carry \\
\end{tabularx}
\end{patternbox}

\exlegend
\begin{enumerate}[leftmargin=*, label=\textbf{E4.\arabic*}, itemsep=6pt]
    \item \exeasy\ \pconcept \exprob{Convert 1D Array Into 2D Array}{LeetCode 2022}{https://leetcode.com/problems/convert-1d-array-into-2d-array/}\ \lc{Easy}. Fill an $m \times n$ matrix row by row with the elements of \texttt{original} (length up to $5\cdot10^4$, $m, n \le 4\cdot10^4$); return an empty matrix if this is impossible.
    \item \exeasy\ \pconcept \cf{800}\ \exprob{Upscaling}{Codeforces 1950B}{https://codeforces.com/problemset/problem/1950/B}. Print a $2n \times 2n$ checkerboard made of $2 \times 2$ blocks, starting with \texttt{\#\#} in the top-left corner ($n \le 20$).
    \item \exeasy\ \precent \exprob{Snake in Matrix}{LeetCode 3248}{https://leetcode.com/problems/snake-in-matrix/}\ \lc{Easy}. In an $n \times n$ grid whose cell $(i,j)$ has id $i\cdot n + j$, a snake starts at cell 0 and follows up to 100 commands \texttt{UP}, \texttt{RIGHT}, \texttt{DOWN}, \texttt{LEFT} (it never leaves the grid). Return the id of its final cell.
    \item \exeasy\ \ptricky \cf{800}\ \exprob{Good Kid}{Codeforces 1873B}{https://codeforces.com/problemset/problem/1873/B}. Given $n \le 9$ digits, add 1 to exactly one of them so that the product of all digits is as large as possible; print that product.
    \item \exeasy\ \precent \exprob{Count Tested Devices After Test Operations}{LeetCode 2960}{https://leetcode.com/problems/count-tested-devices-after-test-operations/}\ \lc{Easy}. Devices $0..n-1$ are visited in order. If the battery of device $i$ is positive, it is tested and the battery of every later device drops by 1 (never below 0). Count the tested devices.
    \item \exeasy\ \pclassic \cf{800}\ \exprob{Sum of Round Numbers}{Codeforces 1352A}{https://codeforces.com/problemset/problem/1352/A}. A number is round if all its digits except the first are zeros. Write $n \le 10^4$ as a sum of as few round numbers as possible.
    \item \exeasy\ \ptricky \exprob{Pass the Pillow}{LeetCode 2582}{https://leetcode.com/problems/pass-the-pillow/}\ \lc{Easy}. People $1..n$ stand in a line; person 1 holds a pillow. Every second it is passed to the next person; at either end of the line the direction reverses. Who holds it after \texttt{time} seconds? Solve it for $n, \texttt{time} \le 1000$, then for $\texttt{time} \le 10^{18}$.
\end{enumerate}

\subsubsection*{Hints}
\begin{enumerate}[leftmargin=*, label=\textbf{E4.\arabic*}]
    \item Element $i$ of \texttt{original} goes to which row and column? When is the matrix impossible?
    \item Which $2 \times 2$ block contains cell $(r, c)$? The block colours alternate like a chessboard.
    \item Keep the row and the column, not the id; convert at the end.
    \item Which digit should receive the $+1$? Think about a 0 among the digits, and about the ratio $(d+1)/d$.
    \item You never need to modify the array: after $t$ tests, how much has every later battery dropped?
    \item Every non-zero digit contributes one round number; which one?
    \item The motion is periodic. How long is one round trip? Where is the pillow at each time of the first period?
\end{enumerate}

\subsubsection*{Answers}

\begin{solutionbox}[E4.1 Convert 1D Array Into 2D Array]
Row-major filling puts \texttt{original[i]} at row $\lfloor i/n \rfloor$, column $i \bmod n$ ($n$ = number of \emph{columns}). It is possible iff $m \cdot n$ equals the length; compute $m \cdot n$ in \texttt{long long}, since $4\cdot10^4 \cdot 4\cdot10^4 = 1.6\cdot10^9$ is close to the \texttt{int} limit.
\begin{lstlisting}[style=cpp]
#include <bits/stdc++.h>
using namespace std;

class Solution {
public:
    vector<vector<int>> construct2DArray(vector<int>& original, int m, int n) {
        if ((long long)m * n != (long long)original.size()) return {};
        vector<vector<int>> res(m, vector<int>(n));
        for (int i = 0; i < (int)original.size(); i++)
            res[i / n][i % n] = original[i];        // 1D index -> (row, column)
        return res;
    }
};

int main() {
    vector<int> a = {1, 2, 3, 4}, b = {1, 2};
    vector<vector<int>> r = Solution().construct2DArray(a, 2, 2);
    cout << r[0][0] << r[0][1] << " " << r[1][0] << r[1][1] << " | "
         << Solution().construct2DArray(b, 1, 1).size() << "\n";   // 12 34 | 0
}
\end{lstlisting}
Official samples: $[1,2,3,4]$, $2 \times 2$ $\to$ $[[1,2],[3,4]]$; $[1,2,3]$, $1 \times 3 \to [[1,2,3]]$; $[1,2]$, $1 \times 1 \to [\,]$.
\end{solutionbox}

\begin{solutionbox}[E4.2 Upscaling]
Cell $(r, c)$ lies in block $(\lfloor r/2 \rfloor, \lfloor c/2 \rfloor)$, and blocks alternate colours like a chessboard: print \texttt{\#} iff $\lfloor r/2 \rfloor + \lfloor c/2 \rfloor$ is even.
\begin{center}
\begin{tikzpicture}[x=0.55cm, y=0.55cm]
    \foreach \r in {0,...,5}{ \foreach \c in {0,...,5}{
        \pgfmathtruncatemacro{\b}{mod(int(\r/2) + int(\c/2), 2)}
        \ifnum\b=0 \fill[primaryblue!60] (\c,-\r) rectangle ++(1,-1); \else \draw[lightgray] (\c,-\r) rectangle ++(1,-1); \fi
    }}
    \draw[thick, accentorange] (2,-2) rectangle ++(2,-2);
    \node[note, anchor=west, text width=5.2cm, align=left] at (6.8,-3) {$n = 3$: the orange block is block $(1,1)$; $1 + 1$ even $\Rightarrow$ \texttt{\#}. Cell $(3,2)$: $\lfloor 3/2 \rfloor + \lfloor 2/2 \rfloor = 2$, even.};
\end{tikzpicture}
\end{center}
\begin{lstlisting}[style=cpp]
#include <bits/stdc++.h>
using namespace std;

int main() {
    int t;
    cin >> t;
    while (t--) {
        int n;
        cin >> n;
        for (int r = 0; r < 2 * n; r++) {
            for (int c = 0; c < 2 * n; c++)
                cout << (((r / 2 + c / 2) % 2 == 0) ? '#' : '.');
            cout << "\n";
        }
    }
}
\end{lstlisting}
The output for the official sample ($n = 1, 2, 3, 4$) matches character by character.
\end{solutionbox}

\begin{solutionbox}[E4.3 Snake in Matrix]
Keep $(i, j)$, start at $(0,0)$; \texttt{UP}/\texttt{DOWN} change $i$, \texttt{LEFT}/\texttt{RIGHT} change $j$; return $i \cdot n + j$. The guarantee ``never leaves the grid'' removes the border checks. Samples: $n = 2$, \texttt{RIGHT, DOWN} $\to (1,1) \to 3$; $n = 3$, \texttt{DOWN, RIGHT, UP} $\to (0,1) \to 1$. Simulating on the id directly ($\pm 1$, $\pm n$) works too, but only because the snake never wraps around a row end --- the $(i, j)$ form is the safe habit.
\end{solutionbox}

\begin{solutionbox}[E4.4 Good Kid]
Adding 1 to digit $d$ multiplies the product by $(d+1)/d$, which is largest for the \textbf{smallest} digit. \textbf{The trap} is the digit 0: if there is a 0, the product is 0 unless that very 0 becomes 1, so ``increase the smallest digit'' is still right (and it is \emph{the} only choice when there is exactly one 0). With two zeros the answer is 0 whatever we do.
\begin{lstlisting}[style=cpp]
#include <bits/stdc++.h>
using namespace std;

int main() {
    int t;
    cin >> t;
    while (t--) {
        int n;
        cin >> n;
        vector<int> a(n);
        for (auto& x : a) cin >> x;
        *min_element(a.begin(), a.end()) += 1;      // the smallest digit gains the most
        long long p = 1;                            // 10^9 > int range in general
        for (int x : a) p *= x;
        cout << p << "\n";
    }
}
\end{lstlisting}
Samples: $[2,2,1,2] \to 16$, $[0,1,2] \to 2$, $[4,3,2,3,4] \to 432$, nine 9s $\to 10 \cdot 9^8 = 430467210$. Checked against trying every digit on random inputs with digits $0..9$.
\end{solutionbox}

\begin{solutionbox}[E4.5 Count Tested Devices After Test Operations]
Literal simulation is $\bigO{n^2}$ (fine for $n \le 100$). But after $t$ tests every later battery has dropped by exactly $t$ (clamped at 0), so device $i$ is tested iff $b_i - t > 0$. One counter replaces the inner loop: $\bigO{n}$.
\begin{lstlisting}[style=cpp]
#include <bits/stdc++.h>
using namespace std;

class Solution {
public:
    int countTestedDevices(vector<int>& b) {
        int tested = 0;
        for (int x : b)
            if (x - tested > 0) tested++;   // battery after `tested` decrements
        return tested;
    }
};

int main() {
    vector<int> a = {1, 1, 2, 1, 3}, c = {0, 1, 2};
    cout << Solution().countTestedDevices(a) << " "
         << Solution().countTestedDevices(c) << "\n";   // 3 2
}
\end{lstlisting}
Sample $[1,1,2,1,3]$: device 0 ($1 - 0 > 0$, tested, $t = 1$), device 1 ($1 - 1 = 0$, no), device 2 ($2 - 1 > 0$, $t = 2$), device 3 ($1 - 2 < 0$), device 4 ($3 - 2 > 0$, $t = 3$) $\Rightarrow 3$. The clamp ``never below 0'' does not matter, because a clamped battery would only be compared after more decrements. Checked against the literal $\bigO{n^2}$ simulation.
\end{solutionbox}

\begin{solutionbox}[E4.6 Sum of Round Numbers]
Read the digits of $n$ from the least significant end with $n \bmod 10$ and $n / 10$, keeping the place value $p = 1, 10, 100, \ldots$; every non-zero digit $d$ gives the round number $d \cdot p$. Fewer terms are impossible: adding one round number to any number raises the count of non-zero digits by at most one (a carry only turns 9s into 0s and then raises a single digit), so $k$ round numbers have a sum with at most $k$ non-zero digits. The number of non-zero digits of $n$ is therefore a lower bound, and the construction reaches it.
\begin{lstlisting}[style=cpp]
#include <bits/stdc++.h>
using namespace std;

int main() {
    int t;
    cin >> t;
    while (t--) {
        int n;
        cin >> n;
        vector<int> parts;
        for (int p = 1; n > 0; p *= 10, n /= 10)
            if (n % 10 != 0) parts.push_back(n % 10 * p);   // digit times its place value
        cout << parts.size() << "\n";
        for (int x : parts) cout << x << " ";
        cout << "\n";
    }
}
\end{lstlisting}
Sample: $5009 \to 9 + 5000$, $9876 \to 6 + 70 + 800 + 9000$, $10000 \to 10000$ (the judge accepts any order).
\end{solutionbox}

\begin{solutionbox}[E4.7 Pass the Pillow]
Simulation: keep (position, direction) and reverse the direction at $1$ and $n$; $\bigO{\texttt{time}}$. \textbf{The trap} is the period: a round trip $1 \to n \to 1$ takes $2(n-1)$ seconds, not $2n$ (the turn-around person does not hold the pillow twice). With $t = \texttt{time} \bmod 2(n-1)$ the holder is $1 + t$ if $t < n$, otherwise $2n - 1 - t$.
\begin{center}
\begin{tikzpicture}[x=0.8cm, y=1cm]
    \foreach \t/\p in {0/1, 1/2, 2/3, 3/4, 4/3, 5/2, 6/1, 7/2}{
        \node[idx] at (\t,1.25) {\t};
        \ifnum\t<6 \node[cell, minimum width=7mm] (q\t) at (\t,0.6) {\p}; \else \node[aout, minimum width=7mm] (q\t) at (\t,0.6) {\p}; \fi
    }
    \node[rowlabel, font=\scriptsize\color{darkgray}] at (-0.55,1.25) {time};
    \node[rowlabel] at (-0.55,0.6) {holder};
    \draw[decorate, decoration={brace, mirror, amplitude=4pt}, thick, memoorange] (-0.4,0.15) -- (5.4,0.15) node[midway, below=4pt, font=\scriptsize\bfseries, text=memoorange] {period $2(n-1) = 6$ for $n = 4$};
\end{tikzpicture}
\end{center}
\begin{lstlisting}[style=cpp]
#include <bits/stdc++.h>
using namespace std;

long long passThePillow(long long n, long long time) {
    long long t = time % (2 * (n - 1));        // one round trip takes 2(n-1) seconds
    return t < n ? 1 + t : 2 * n - 1 - t;
}

int main() {
    cout << passThePillow(4, 5) << " " << passThePillow(3, 2) << " "
         << passThePillow(1000, 1000000000000000000LL) << "\n";   // 2 3 999
}
\end{lstlisting}
Samples: $n = 4$, \texttt{time} $= 5$ $\to 2$; $n = 3$, \texttt{time} $= 2 \to 3$; and $n = 1000$, $\texttt{time} = 10^{18}$: $t = 10^{18} \bmod 1998 = 1000 \ge n$, holder $1999 - 1000 = 999$. The formula agrees with the step-by-step simulation for all $n \le 7$, $\texttt{time} \le 40$.
\end{solutionbox}
